Research on Generation Z and workplace ethics has mostly asked whether this cohort holds stronger ethical values than earlier ones, a premise that reviews of generational research do little to support. This conceptual paper asks a different question: when employees conclude that their employer has broken an ethical commitment, where do they take the complaint? Grounded in psychological contract theory, especially its account of ideological currency, and in the exit-voice-loyalty framework, the paper develops a venue congruence model. It introduces promise publicity, the extent to which an employee learned of the employer's ethical obligations through communication addressed to the public, and proposes that breaches of public ethical promises are more likely to be felt as violations and voiced publicly. Perceived internal voice safety is proposed as the boundary that decides whether violation is voiced inside, taken public, or turned into exit or silence. Rival propositions separate a cohort explanation of the Generation Z pattern from a career-stage explanation. The paper reports no data and outlines designs for testing its five propositions.
